\section{Теоретические основы многосторонних конфиденциальных вычислений}
\label{sec:theoretical-foundations}

\subsection{Определение и общая характеристика МКВ}

Формально под многосторонним конфиденциальным вычислением понимается протокол \(\pi\), который принимает от каждого участника \(P_{i}\) его секретный аргумент \(x_{i}\in D\) и совместно с другими участниками вычисляет значение некоторой известной функции \(F : D^{n}\to R\). Протокол должен удовлетворять двум базовым свойствам: \textit{корректность} --- все честные участники получают верный результат \(y=F(x_{1},\ldots,x_{n})\); \textit{конфиденциальность} --- никакая информация о входах честных сторон кроме того, что следует из результата вычислений, не раскрывается другим участникам. Это означает, что протокол обеспечивает информативный эквивалент присутствия доверенной стороны \(T\), которая принимает от всех сторон их входы и возвращает только результат.

Математическая постановка задачи включает описание идеального и реального мира. В идеальном мире существует абсолютно доверенная функция \(T\), которой стороны передают свои входы \(x_{i}\). Функция вычисляет \(F(x_{1},\ldots,x_{n})\) и отправляет результат всем сторонам; в этом мире угрозы отсутствуют, а конфиденциальность и корректность гарантированы. В реальном мире такой функции нет, и стороны взаимодействуют между собой по протоколу \(\pi\). Цель МКВ --- построить протокол, который делает неразличимыми (с точки зрения вычислительной ограниченности злоумышленника) распределения выходов и представлений (всего, что видит участник или злоумышленник во время выполнения протокола) сторон в идеальном и реальном мирах. В научной литературе это называют реально-идеальной парадигмой; она задаёт формальную основу для доказательств безопасности протоколов.

Помимо корректности и конфиденциальности в МКВ выделяют дополнительные свойства: устойчивость (robustness), прекращаемость и справедливость (fairness). Устойчивость означает, что протокол сохраняет корректность вычисления при отказе части участников или при ошибках каналов связи. Справедливость гарантирует, что если одна честная сторона получает результат, то его получают и остальные честные стороны; однако это свойство достижимо только в ограниченных моделях, например при использовании доверенного компонента или экономических стимулов. Поэтому большинство практических протоколов ограничивается гарантиями конфиденциальности и корректности и допускает безопасное досрочное прекращение вычисления при обнаружении атаки.

\subsection{Исторические предпосылки и развитие МКВ}

Истоки МКВ связаны с ментальным покером --- системой криптографических задач 1970-х гг., позволявшей проводить карточные игры без доверенного дилера. Теоретические основы области заложил Эндрю Яо: в 1982 г. он сформулировал задачу миллионеров, а в 1986 г. предложил общий подход к безопасному двустороннему вычислению; термин искаженная схема (garbled circuit) закрепился в литературе позднее \cite{14, 15, 16}. В конце 1980-х гг. появились общие многосторонние протоколы и формальные модели безопасности, показавшие возможность защищённого вычисления произвольных функций \cite{06, 08}. С 2010-х гг. исследования сосредоточены на эффективности и практическом внедрении МКВ, а создание в 2019 г. MPC Alliance отразило переход технологии к отраслевому применению \cite{19}.

\subsection{Формальная постановка задачи}

Пусть имеется конечное множество участников \(P_{1},\ldots,P_{n}\). Каждый участник \(P_{i}\) обладает секретным входом \(x_{i}\). Все участники заинтересованы в вычислении результата \(y=F(x_{1},\ldots,x_{n})\), где \(F\) --- детерминированная или вероятностная функция, известная заранее. В идеальной модели с доверенным посредником \(T\) каждый \(P_{i}\) отправляет \(x_{i}\) посреднику, который выдаёт всем результат \(y\). В реальном протоколе отсутствует доверенный посредник, поэтому участники отправляют сообщения друг другу по заранее определённым алгоритмам. Цель протокола \(\pi\) --- обеспечить, чтобы распределения выходных данных и информационных следов сторон в реальном мире были вычислительно неотличимы от таковых в идеальной модели. Для получестного (semi-\allowbreak honest) нарушителя формальное определение описывается через два распределения случайных переменных. Распределение результатов и наблюдений при реальном выполнении протокола обозначается:

\begin{equation}
\mathrm{Real}_{\pi}(\kappa, C; x_{1},\ldots,x_{n})
\end{equation}

где \(\mathrm{Real}_{\pi}\) --- результат выполнения протокола \(\pi\) с параметром безопасности \(\kappa\) и множеством коррумпированных сторон \(C\); распределение включает виды всех коррумпированных сторон и выходы всех участников. В идеальной модели с доверенной функцией и симулятором соответствующее распределение обозначается:

\begin{equation}
\mathrm{Ideal}_{F,\mathrm{Sim}}(\kappa, C; x_{1},\ldots,x_{n})
\end{equation}

где \(\mathrm{Ideal}_{F,\mathrm{Sim}}\) --- распределение в идеальной модели: доверенная функция вычисляет результат, а симулятор \(\mathrm{Sim}\) восстанавливает виды коррумпированных сторон на основе входов и выходов. Протокол безопасно реализует функцию \(F\), если для любого множества коррумпированных сторон эти распределения неразличимы.

Приведём пример: пусть три участника хотят вычислить максимум своих зарплат, не раскрывая их. Они хотят вычислить функцию \(F(x,y,z)=\max(x,y,z)\). В идеальной модели каждый отправляет свой вход доверенному посреднику, который возвращает результат. В реальном мире можно построить протокол на основе разделения секрета: каждый участник делит свой вход на доли и использует операции над зашифрованными битами. Другой вариант --- применить протокол искаженных схем: составить схему сравнения и обменяться зашифрованными метками. Таким образом достигается корректное вычисление \(\max(x,y,z)\) без раскрытия \(x,y,z\) каждому отдельному участнику.

\subsection{Модель вычислений и модели нарушителя}

Модель МКВ задаёт допущения о каналах связи и типах нарушителей. Протоколы МКВ предполагают, что между каждой парой участников имеется аутентифицированный и конфиденциальный канал передачи, что позволяет исключить перехват и модификацию сообщений внешним злоумышленником. Внутренний злоумышленник может нарушать правила протокола, поэтому различают два основных типа:

\begin{itemize}
\item Получестный (semi-\allowbreak honest) нарушитель, или (иногда) честный, но любопытный. Он корректно исполняет инструкции протокола, но пытается по полученным сообщениям восстановить чужие данные. Формально такой нарушитель может объединять представления коррумпированных сторон, но не изменяет ход вычислений. Безопасность протокола обеспечивается наличием симулятора, который в идеальном мире воспроизводит представление нарушителя, что гарантирует отсутствие утечки информации.

\item Злонамеренный (malicious) нарушитель может отклоняться от протокола, отправлять поддельные сообщения, прекращать участие или попытаться изменить вывод честных сторон. Обеспечение безопасности против такого нарушителя требует усиленных механизмов контроля, таких как доказательства с нулевым разглашением, проверка корректности операций и алгоритм "Дели и выбирай" (cut-and-choose). Важно, что многие практические протоколы сначала строятся и анализируются для получестной модели, а затем адаптируются для защиты от злонамеренного нарушителя с использованием протоколов GMW или SPDZ (Shamir-based Protocol for Distributed ZK-proofs) \cite{06}.
\end{itemize}

Кроме того, в литературе выделяют модели сетевых задержек (синхронная, полусинхронная, асинхронная), предположения о доле честных участников (например, честное большинство: если \(t\) --- число коррумпированных участников, то \(t<n/2\); такой порог используется, в частности, в некоторых информационно-\allowbreak теоретических протоколах против получестного нарушителя, тогда как для активного нарушителя допустимый порог может быть строже) и возможность наличия дополнительных примитивов (доступ к общему радио-каналу или «чёрной доске») \cite{06}, \cite{08}. Эти параметры влияют на возможность достижения критерия справедливости: как показано в ряде работ, для общего числа участников выполнить вычисление с гарантированным выводом без доверенной стороны невозможно, если злоумышленник может прерывать коммуникацию, что известно как теорема Клива \cite{08} --- большинство протоколов реализует только безопасность с возможностью досрочного прекращения.

\subsection{Основные свойства МКВ}

Свойства протокола МКВ можно классифицировать следующим образом:

\begin{enumerate}
\item Конфиденциальность. Никакая информация о входах честных участников, кроме того, что следует из результата, не может быть выведена злоумышленником. Этого достигают за счёт использования разделения секрета, зашифрованных схем, гомоморфного шифрования и других криптографических примитивов.

\item Корректность. Результат, получаемый в ходе протокола, совпадает с результатом идеальной функции \(F\); для получестной модели корректность следует напрямую из определения безопасности, для модели со злонамеренным нарушителем требуются дополнительные проверки.

\item Устойчивость и отказоустойчивость. Протокол способен завершить вычисление при отказе части участников. Многие протоколы на основе разделения секрета допускают восстановление результата при наличии определённого числа честных участников (t-out-of-n).

\item Справедливость. Если хотя бы один честный участник получает вывод, то все честные участники также должны его получить. Достижение справедливости требует специальных механизмов, так как в противном случае злоумышленник может получить результат и прекратить выполнение протокола. В научной литературе гарантированный вывод рассматривается как более сильное свойство, чем справедливость: он влечёт справедливость, однако обратное в общем случае неверно \cite{24}. Его достижение возможно при честном большинстве или с применением дополнительных технологий, таких как защищённые окружения (TEE) или экономические стимулы.
\end{enumerate}

\subsection{Примеры и математические иллюстрации}
\subsubsection{Секретное суммирование}
Пусть \(n\) участников хотят вычислить сумму своих чисел \(s=\sum_{i=1}^{n}x_{i}\). В простой схеме разделения секрета каждый участник выбирает случайные числа \(r_{i}^{(j)}\) для \(j=1\ldots n-1\) так, чтобы сумма этих случайных долей была равна \(x_{i}\). Затем участник передаёт \(r_{i}^{(j)}\) j-му соседу, а оставшуюся долю оставляет себе. Получив все доли, участники локально суммируют полученные значения, и каждый получает долю суммы. Суммируя доли, можно восстановить общий результат, но отдельные входы остаются скрытыми. Такая схема является частным случаем общего механизма разделения секрета, например схемы Шамира, где секрет \(s\) представляет собой значение полинома в точке 0, а доли --- значения этого полинома в точках 1,2, ...

\subsubsection{Схема Шамира}

Рассмотрим схему разделения секрета \((t,n)\) по Шамиру. Дилер выбирает случайный полином степени \(t-1\) над полем \(\mathbb{F}\):

\begin{equation}
f(x)=s+\alpha_{1}x+\alpha_{2}x^{2}+\cdots+\alpha_{t-1}x^{t-1},
\end{equation}

где коэффициенты \(\alpha_{i}\) выбираются случайно. Долю участника \(P_{i}\) определяют как \(s_{i}=f(i)\). Любые \(t\) долей позволяют восстановить секрет с помощью интерполяции Лагранжа (см. формулу~\eqref{eq:lagrange-interpolation}), а любой набор из менее чем \(t\) долей не даёт информации о секрете (перфектная конфиденциальность). Эта идея широко используется в протоколах BGW и SPDZ; параметры \(t\) и \(n\) определяют порог устойчивости.

\subsubsection{Искаженные схемы Яо }

Другим механизмом является модель искажённых схем Яо. При её использовании одна сторона («генератор») представляет функцию \(F\) в виде булевой схемы и сопоставляет каждому проводу схемы две случайные метки, соответствующие логическим значениям 0 и 1. Схема преобразуется в «зашифрованную» форму: таблицы заменены на зашифрованные значения. Вторая сторона («вычислитель») с помощью протокола слепой передачи (oblivious transfer) получает метки, соответствующие значениям собственных входных битов, и вычисляет зашифрованную схему, не узнавая открытых значений входов генератора. Полученную выходную метку вычислитель преобразует в результат с помощью таблицы декодирования, предоставленной генератором. Такой подход даёт возможность строить двухсторонние протоколы для произвольных булевых функций, но требует значительных ресурсов памяти и времени.

\subsection{Вывод}

В данном разделе были рассмотрены теоретические основы МКВ: история развития, формальная модель вычислений, модели нарушителя и свойства безопасности; приведены примеры секретного суммирования и схем Яо. Во второй главе будут рассмотрены криптографические основы МКВ: разделение секрета, искаженные схемы, гомоморфное шифрование и протокол слепой передачи.